\documentclass[10pt,a4paper]{amsart}
\usepackage[margin=19mm]{geometry}
\usepackage{amsmath,amssymb,mathtools}
\usepackage[scr=rsfs,cal=euler]{mathalfa}
\usepackage{xfrac}
\usepackage[unicode,hidelinks,bookmarks=false]{hyperref}
\newcommand{\ie}{i.e.\ }
\newcommand{\cf}{cf.\ }
\newcommand{\resp}{respectively\ }
\newcommand{\Subsection}{\subsection}
\providecommand{\nobreakdash}{\nobreak\hbox{-}}
\setlength{\emergencystretch}{2em}
\multlinegap0pt
\usepackage{amssymb}
\usepackage[shortlabels]{enumitem}
\usepackage{bm}
\usepackage{xcolor}
\newtheorem{lemma}{Lemma}
\title{Orthomorphism Polynomials of degree $7$ over finite fields}
\author{Bhitali Kousik, Dhiren Kumar Basnet}
\date{Source: arXiv:2507.10316v2 (CC BY 4.0)}
\begin{document}
\maketitle

\noindent\emph{Source and licence.} Orthomorphism Polynomials of degree $7$ over finite fields. Version arXiv:2507.10316v2 (CC BY 4.0), distributed by arXiv under the Creative Commons Attribution 4.0 International licence, http://creativecommons.org/licenses/by/4.0/, as recorded in the arXiv licence metadata for this exact version and shown on the abstract page https://arxiv.org/abs/2507.10316v2. Full licence notice: \texttt{licenses/arxiv-2507.10316-CC-BY-4.0.txt}.\par\smallskip

\noindent\textit{Editorial scope.} Counted unit: the article's lemma l6 (source0.tex lines 168--170). The article's complete original proof of it is delivered with it (source0.tex lines 171--174, one \texttt{proof} environment of 4 lines). No mathematics is written, repaired or re-derived: the statement and the proof are the article's own lines, in the article's own notation, and no retained line is substituted or re-flowed.  No further source-local context is needed: every cross-reference of every retained line resolves inside the two delivered spans.  Nothing was omitted from any retained span, so the package carries no omission record.  No retained line contains a citation, so the wrapper carries no bibliography and no bibliography entry.  No retained line contains a figure environment, an image-inclusion command or drawing code, so no image is delivered and the package declares no asset dependency.  Editorial formatting only, in addition: an AMS \texttt{amsart} preamble that restates the article's own declarations and macros used by the retained lines, namely the environments \texttt{lemma} on separate counters, together with the standard packages those lines need.  The retained environments print in this document as Lemma 1 in order of appearance, whereas the article prints the same items with the numbering of its own full document.  The complete native source of the article as distributed in the arXiv e-print for this exact version is bundled unchanged as \texttt{sources/181548/source0.tex}.
\begin{lemma}\label{l6}
    Let $F(x)= x^7+\sum_{i=1}^{5}f_ix^i$ be a polynomial over $\mathbb{F}_q.$ Define $g(x)=\alpha F(\beta x)$ and $h(x)=\alpha_1 F(\beta_1 x),$ where $\alpha,~\beta,~\alpha_1,~\beta_1 \in \mathbb{F}_q^*.$~If $g( x+\gamma)+\delta= h(x+\gamma_1)+\delta_1$ for all $x\in \mathbb{F}_q$ with $\gamma,~\delta,~\gamma_1,~\delta_1 \in \mathbb{F}_q,$ then we must have $\gamma=\gamma_1$ and $\delta=\delta_1.$ 
\end{lemma}

\begin{proof}
Equating the coefficients of all powers of $x$ in the given equality with some simplification yields the required result. Details are left to the reader.

\end{proof}

\end{document}
